\chapter{sotFS transaccional: crash refinement}
\label{cap:sotfs-tx}

\section{Motivación}

Un filesystem que vive en disco tiene que sobrevivir
\emph{crashes}: el power-cut, el kernel panic, el hot-yank del
SSD. La pregunta no es \emph{«¿se rompe?»} sino \emph{«¿en qué
estado queda?»}. Hay tres respuestas posibles, en orden creciente
de garantías:

\begin{description}
  \item[\emph{Inconsistent}.] El FS queda corrupto; \texttt{fsck}
    intenta repararlo con heurísticas. Bueno hasta el primer
    archivo crítico que no se recuperó.
  \item[\emph{Eventually consistent}.] Tras el reboot, una
    secuencia de pasos automáticos (replay del journal,
    reconciliación) llega a un estado consistente \emph{eventualmente}.
    ext4 con \texttt{data=journal} entra acá.
  \item[\emph{Always consistent}.] El FS está \emph{en todo
    momento} en un estado consistente. Tras crash, la primera
    lectura ya retorna datos correctos. Es lo que sotFS apunta.
\end{description}

La técnica para conseguir la tercera son las \emph{transacciones
del filesystem} (GTXN), construidas encima de las primitivas
del cap.~\ref{cap:transacciones} pero con la dimensión adicional
de \emph{durabilidad}: una transacción committed \emph{persiste
al disco antes de retornar control al usuario}. La verificación
formal usa \emph{crash refinement}: el spec abstracto
\texttt{sotfs\_transactions.tla} no tiene crashes; el spec
concreto \texttt{sotfs\_crash.tla} sí, y se prueba que el
concreto \emph{refina} el abstracto (toda traza concreta es
observable como traza abstracta).

\section{El protocolo GTXN}

\begin{algorithm}
\caption{\texttt{gtxn\_commit(tx)} --- WAL + fsync}
\begin{algorithmic}[1]
\State \(\triangleright\) Fase 1: WAL escrito en disco
\State $\texttt{wal\_record} \gets \texttt{serialize(tx.changes)}$
\State \texttt{disk.write(wal\_log, wal\_record)}
\State \texttt{disk.fsync(wal\_log)}\Comment{barrier durabilidad}
\State \(\triangleright\) Fase 2: marcar WAL committed (1-byte atomic)
\State \texttt{disk.write(wal\_log, COMMIT\_MARKER)}
\State \texttt{disk.fsync(wal\_log)}\Comment{commit point}
\State \(\triangleright\) Fase 3: aplicar al graph store
\For{cambio \texttt{c} en \texttt{tx.changes}}
  \State \texttt{disk.write(graph\_store, c)}
\EndFor
\State \texttt{disk.fsync(graph\_store)}
\State \(\triangleright\) Fase 4: GC del WAL
\State \texttt{wal.truncate(tx)}
\end{algorithmic}
\end{algorithm}

El \emph{commit point} es el segundo \texttt{fsync} (línea 5).
\emph{Antes} del commit point, el crash deja WAL parcial sin
marker → recovery descarta las operaciones (vuelve al estado
pre-tx). \emph{Después} del commit point, el crash deja WAL
con marker → recovery re-aplica las operaciones (continúa al
estado post-tx). Nunca un estado intermedio.

\section{Crash refinement en TLA}

El spec abstracto \texttt{sotfs\_transactions.tla} modela GTXN
sin crash. El spec concreto \texttt{sotfs\_crash.tla} agrega:

\begin{lstlisting}[language=tla, frame=single, basicstyle=\ttfamily\scriptsize]
\* Crash puede ocurrir en cualquier momento.
Crash ==
    /\ ~crashed
    /\ crashed' = TRUE
    /\ memState' = NULL                  \* memoria volátil borrada
    /\ UNCHANGED <<diskState, walState, gtxnState, ...>>

\* Recovery: replay del WAL desde el último commit marker.
Recover ==
    /\ crashed
    /\ memState' = ReplayWAL(walState, diskState)
    /\ crashed' = FALSE
    /\ UNCHANGED <<diskState, walState, ...>>
\end{lstlisting}

\subsection*{Cómo se lee un \texttt{.tla} (cuarta lectura)}

Después de cap.~\ref{cap:scheduler}, \ref{cap:ipc-endpoint},
\ref{cap:transacciones}, \ref{cap:tla-pipeline}, una
construcción nueva:

\begin{description}
  \item[\texttt{ReplayWAL(\dots)}] una función definida en el
    \texttt{LET} del spec, que toma el estado del WAL y del disco
    y computa el estado de memoria \emph{como si reaplicáramos
    todas las operaciones committed}. La función es \emph{pura}
    (no muta), ideal para razonar sin acción.
\end{description}

La propiedad de \emph{refinement}:

\begin{lstlisting}[language=tla, frame=single, basicstyle=\ttfamily\scriptsize]
RefinementMapping ==
    LET
        ConcreteState == <<diskState, walState, memState>>
        AbstractState ==
          IF crashed THEN
              \* En estado crash, mapeamos al estado abstracto
              \* "como si la última committed tx fuera la actual"
              ReplayWAL(walState, diskState)
          ELSE
              memState
    IN
        \A op \in TxOps :
            \* Toda operación abstracta es una operación concreta
            ConcreteOp(op) => AbstractOp(op)

THEOREM Spec_Concrete => Spec_Abstract
\end{lstlisting}

\section{Las propiedades verificadas}

\begin{description}
  \item[\textbf{No partial application.}] Tras Recover,
    \emph{toda} operación de una GTXN está aplicada o
    \emph{ninguna}. Nunca «se aplicaron 3 de 5».
  \item[\textbf{Durability.}] Si \texttt{gtxn\_commit} retornó
    \texttt{Ok}, la transacción sobrevive cualquier crash
    posterior.
  \item[\textbf{Recovery consistency.}] Tras Recover, el estado
    de memoria satisface todas las invariantes del cap.~\ref{cap:sotfs-modelo}
    (DirIntegrity, ExtentIntegrity, NameUniqueness).
  \item[\textbf{WAL integrity.}] Una vez escrito el commit
    marker, el contenido del WAL no se modifica hasta el GC; el
    GC sólo borra entradas \emph{ya aplicadas}.
\end{description}

\section{El protocolo de fsync}

La barrera de durabilidad es \texttt{fsync}: pide al SSD que
flushee sus cachés volátiles a flash NAND. Sin \texttt{fsync},
una escritura puede quedarse en el cache del controlador y
perderse en un power-cut. El protocolo:

\begin{lstlisting}[frame=single, basicstyle=\ttfamily\scriptsize]
disk.write(buf)         # bytes a kernel page cache
disk.fsync()            # kernel page cache → driver buffer
                        # driver buffer → SSD controller
                        # SSD controller → NAND
                        # ack al kernel
\end{lstlisting}

\begin{invariante}[Persistencia tras fsync]
\label{inv:fsync-persist}
Tras \texttt{disk.fsync(loc)} retornar exitosamente, los bytes
escritos antes de la llamada a \texttt{loc} están persistidos en
NAND y sobrevivirán cualquier power-cut posterior.
\end{invariante}

La invariante depende del SSD respetar el comando \texttt{FUA}
(\emph{Force Unit Access}). SSDs consumer-grade malos lo
ignoran; SSDs enterprise-grade lo respetan. Verificar esto
runtime con un test de power-cut físico es la única vía; el
modelo asume \texttt{fsync} honesto.

\section{Trampa: write reordering}

\begin{trampa}[\texttt{write A; write B; fsync} no garantiza orden]
Versión vulnerable: el código asume que dos \texttt{write}
seguidos por un \texttt{fsync} se persisten en orden A, B. NO
es así. El SSD puede reordernar internamente; \texttt{fsync}
sólo garantiza que \emph{ambos} están persistidos al retornar,
no el orden.

Si el código depende de que A esté en disco antes que B (\emph{e.g.}~B
es el commit marker que valida A), una crash entre los dos
fsyncs puede dejar B en disco sin A.

Fix: \texttt{write A; fsync; write B; fsync}. Dos
\texttt{fsync} explícitos. Costo: latencia, especialmente en
HDD. Justificable cuando el orden importa (commit marker), no
cuando no (datos paralelos).

\textbf{Lección:} \texttt{fsync} es una barrera de
\emph{durabilidad puntual}, no de \emph{ordering entre
operaciones}. Cada barrier tiene su propio fsync.
\end{trampa}

\section{Trampa: el GC del WAL antes del fsync del graph}

\begin{trampa}[Recovery imposible si WAL ya borrado]
\citaCommit{4775904}. Versión vulnerable: el orden de Fase 3 y
Fase 4 estaba invertido. El WAL se truncaba (Fase 4) \emph{antes}
de que el graph store hiciera \texttt{fsync} (Fase 3 incompleta).

Si crash ocurre después del WAL truncate pero antes del graph
fsync: WAL no tiene la transacción (truncado), graph no la
tiene (no fsynced). \emph{Pérdida silenciosa}.

Fix: el orden estricto Fase 1 → Fase 2 → Fase 3 → Fase 4. Cada
fase termina con su \texttt{fsync} correspondiente. WAL truncate
sólo tras graph fsync exitoso.

\textbf{Lección:} la durabilidad es un pipeline ordenado de
barreras. Una operación de \emph{cleanup} (truncate del WAL)
DEBE preceder de todas las operaciones que dependen de los
datos cleaned-up. Caer en la tentación de hacer cleanup
\emph{en paralelo} para «ahorrar latencia» rompe la cadena de
durabilidad.
\end{trampa}

\section{La invariante de \emph{readers see committed-or-pre-tx}}

\begin{invariante}[Read isolation]
\label{inv:read-iso}
Un \texttt{read} concurrente con un \texttt{gtxn\_commit} en
progreso ve \emph{o} el estado anterior a la transacción \emph{o}
el estado posterior; nunca un estado intermedio. La elección
depende de cuándo se linealiza el read respecto del commit
point.
\end{invariante}

Implementación: los readers usan Tier 0 RCU del
cap.~\ref{cap:transacciones}; toman snapshot de la época
\emph{global} del FS al iniciar la lectura, leen consistentemente
de esa snapshot. Los writers (gtxn) bumpean la época sólo en
su Fase 3 (después del commit point), de modo que los readers
con snapshot vieja siguen viendo el estado pre-tx, los nuevos
ven el post-tx.

\section{Origen del diseño}

WAL viene de bases de datos (ARIES \cite{Mohan1992-ARIES}); su
aplicación a filesystems es de \emph{ext3} y \emph{XFS}. La idea
de \emph{crash refinement} (modelo concreto con crashes
refinando un modelo abstracto sin) es de Yggdrasil
\cite{Sigurbjarnarson2016-Yggdrasil}, donde el filesystem es un
proyecto académico con verificación en Push-Button. La
combinación con un Merkle DAG es contribución original de
\texttt{sotX}; la dificultad es mantener la inmutabilidad del
DAG bajo crashes (los nodos viejos sobreviven, los nuevos
pueden estar parciales).

\section*{Resumen del capítulo}

\begin{itemize}
  \item GTXN protocoliza la mutación durable: WAL → commit
    marker → graph apply → WAL GC, con fsync en cada paso.
  \item El \emph{commit point} es el fsync del marker; antes,
    crash descarta; después, crash re-aplica.
  \item Crash refinement (\citaTLA{sotfs\_crash.tla}) prueba
    que toda traza concreta con crashes es traza abstracta sin
    ellos; las propiedades del cap.~\ref{cap:sotfs-modelo}
    sobreviven.
  \item Read isolation: los readers usan Tier 0 RCU sobre la
    época FS; ven o pre-tx o post-tx, nunca intermedio.
  \item Trampas: write reordering sin fsync intermedio (rompe
    durabilidad ordenada); WAL truncate antes del graph fsync
    (pérdida silenciosa).
\end{itemize}

\section*{Ejercicios}

\begin{ejercicio}[$\star$]
Diga el orden \emph{correcto} de los cuatro fsync en una GTXN.
¿Cuál es el commit point? ¿Qué pasa si crash ocurre entre el
fsync 1 y el fsync 2?
\end{ejercicio}

\begin{ejercicio}[$\star\star$]
Lea la propiedad \emph{Recovery consistency}. ¿Qué propiedades
del cap.~\ref{cap:sotfs-modelo} están en juego? ¿Bastaría
verificar sólo \texttt{DirIntegrity} para garantizar
\emph{consistency}?
\end{ejercicio}

\begin{ejercicio}[$\star\star$]
Un atacante con acceso físico al SSD puede inyectar power-cuts
controlados. ¿Puede forzar pérdida de datos en una GTXN
correctamente implementada? Justifique en términos de la
invariante~\ref{inv:fsync-persist} y los assumptions sobre el
SSD.
\end{ejercicio}

\begin{ejercicio}[$\star\star\star$]
Diseñe un test de \emph{crash injection}: una herramienta que
mata el kernel en momentos aleatorios y verifica que el
filesystem está en un estado consistente tras boot. Pista:
QEMU + monitor command \texttt{quit} + script de assertion.
\end{ejercicio}

\begin{ejercicio}[$\star\star\star$]
Compare las garantías de crash safety de sotFS contra ext4
(con \texttt{data=journal}), btrfs y ZFS. ¿Qué ofrece sotFS que
los otros no? ¿Qué ofrecen los otros que sotFS no?
\end{ejercicio}

% TODO(apendice-E): respuestas a ejercicios 18.1-18.5
